дуже\\
дуже давно,\\
один юнак запитав\\
про всі жахіття світу.\\

невдоволений\\
таким багнистим станом,\\
що здавався людським,\\
він рушив у дорогу до гір.\\

у високому місці\\
він знайшов старий дім\\
дивовижної форми.\\
зацікавившись, подзвонив у дзвін.\\

двері відчинив старець\\
усміхнений,\\
безтурботний,\\
що запросив його сісти.\\

\emph{що ти робиш тут, у високому місці?}\\
запитав старець.\\
\emph{я неспокійний.}\\
відповів юнак, провадячи далі...\\

\emph{світ став потворним.\\
я не бачу сенсу.\\
я не бачу рішення\\
для жахіть світу.}\\

\emph{я в пошуку\\
розуміння,\\
розрізнення,\\
звільнення.}\\

\emph{чи є мудрість,\\
здатна згасити бурі?\\
чи є вчинок,\\
здатний урятувати невинних?}\\

\emph{відповідай мені, прошу.\\
яка твоя думка,\\
які твої роздуми\\
про жахіття світу?}\\

приголомшений,\\
відсапуючись,\\
юнак прихиляє\\
спину до спинки.\\

\emph{я одразу бачу,\\
що ти прийшов здалеку.\\
ти пройшов крізь голод.\\
ти пройшов крізь спрагу.}\\

\emph{занурений у запитання\\
і тривоги,\\
питаєш про мир серця.\\
ти заслуговуєш на гідну відповідь.}\\

\emph{прошу тебе тому\\
трохи терпіння.\\
я дам тобі відповідь.\\
найщирішу відповідь.}\\

старець пішов болісними кроками.\\
сів навпроти стіни\\
на стару подушку\\
і залишився там.
